% Conclusión individual: Joseph
En mi participación me enfoqué en la comunicación mediante sockets TCP y en cómo estructurar el intercambio entre el cliente y el servidor. La implementación usa una conexión que se reutiliza para varios mensajes; cada solicitud viaja como JSON codificado en UTF-8 y termina con un salto de línea para que el servidor pueda separar los mensajes. Al trabajar con este mecanismo entendí que TCP entrega un flujo de bytes y que, por eso, el formato, la delimitación y varias validaciones deben resolverse en la aplicación.

Las pruebas funcionales permitieron comprobar el intercambio de solicitudes y acuses entre los nodos. También aprendí que una conexión puede fallar o agotarse el tiempo de espera, y que esos casos deben distinguirse de una respuesta exitosa. En el benchmark secuencial de localhost, TCP registró menor latencia media que gRPC para los dos tamaños evaluados; aun así, ese resultado corresponde a las implementaciones y condiciones usadas en esta práctica, no a una regla para cualquier sistema.

La comparación me ayudó a entender que elegir un mecanismo de comunicación no depende solo de la velocidad. TCP ofreció control sobre el formato, pero exigió más trabajo manual para validar mensajes y manejar errores. Esta experiencia me permitió relacionar las decisiones de implementación con los resultados medidos y reconocer qué aspectos todavía habría que probar en una red y con otras cargas.